% =====================================================================
%  minggu-02.tex  --  Slide Pertemuan Minggu ke-2
%  Pengantar Kriptografi (MPFA014) -- Prodi Matematika FMIPA UBB
%
%  Topik silabus : Layanan keamanan & serangan
%  Sumber materi : buku_ajar/chapters/02 (Bab 2)
%  Sub-CPMK      : 2.1 dan 2.4
%
%  Kompilasi dari dalam folder ppt/:
%      compile-ppt.bat 02
% =====================================================================

\documentclass[aspectratio=169,11pt]{beamer}
% =====================================================================
%  beamer-preamble.tex  --  prasarana dan gaya slide kuliah
%
%  Mata kuliah : Pengantar Kriptografi (MPFA014)
%  Program     : Program Studi Matematika, FMIPA, Universitas Bale Bandung
%
%  Berkas ini TIDAK memuat \documentclass. Setiap berkas minggu memuatnya
%  sendiri lalu menarik berkas ini:
%
%      \documentclass[aspectratio=169,11pt]{beamer}
%      % =====================================================================
%  beamer-preamble.tex  --  prasarana dan gaya slide kuliah
%
%  Mata kuliah : Pengantar Kriptografi (MPFA014)
%  Program     : Program Studi Matematika, FMIPA, Universitas Bale Bandung
%
%  Berkas ini TIDAK memuat \documentclass. Setiap berkas minggu memuatnya
%  sendiri lalu menarik berkas ini:
%
%      \documentclass[aspectratio=169,11pt]{beamer}
%      % =====================================================================
%  beamer-preamble.tex  --  prasarana dan gaya slide kuliah
%
%  Mata kuliah : Pengantar Kriptografi (MPFA014)
%  Program     : Program Studi Matematika, FMIPA, Universitas Bale Bandung
%
%  Berkas ini TIDAK memuat \documentclass. Setiap berkas minggu memuatnya
%  sendiri lalu menarik berkas ini:
%
%      \documentclass[aspectratio=169,11pt]{beamer}
%      \input{beamer-preamble}
%      \title{...} \subtitle{...} \setWeekInfo{Minggu n}{Topik}
%
%  Berkas minggu hanya berisi \setWeekInfo dan lingkungan frame.
%  Kompilasi dijalankan DARI DALAM folder ppt/ (lihat compile-ppt.bat),
%  karena gambar buku dirujuk dengan jalur relatif ../buku_ajar/.
% =====================================================================

% ---------------------------------------------------------------- bahasa
\usepackage[utf8]{inputenc}
\usepackage[T1]{fontenc}
\usepackage[indonesian]{babel}
\usepackage{lmodern}
\usepackage{microtype}

% ------------------------------------------------------------- matematika
\usepackage{amsmath,amssymb,mathtools}

% ----------------------------------------------------------------- tabel
\usepackage{booktabs}
\usepackage{array}
\usepackage{tabularx}

% --------------------------------------------------------------- grafik
\usepackage{tikz}
\usetikzlibrary{arrows.meta,positioning,shapes.geometric,calc,fit,backgrounds}

% ---------------------------------------------------------- kode program
\usepackage{listings}
\lstdefinestyle{slidecode}{
    basicstyle=\ttfamily\scriptsize,
    keywordstyle=\color{ubbblue}\bfseries,
    commentstyle=\color{ubbgreen}\itshape,
    stringstyle=\color{ubbred},
    numbers=left,
    numberstyle=\tiny\color{gray},
    numbersep=6pt,
    frame=single,
    rulecolor=\color{ubbblue!40},
    backgroundcolor=\color{ubbblue!3},
    breaklines=true,
    showstringspaces=false,
    tabsize=4,
    columns=fullflexible,
    keepspaces=true,
}
\lstset{style=slidecode}

% =====================================================================
%  WARNA INSTITUSI
% =====================================================================
\definecolor{ubbblue}{RGB}{12, 62, 120}      % warna utama
\definecolor{ubbteal}{RGB}{0, 118, 110}      % aksen kedua
\definecolor{ubbgreen}{RGB}{22, 122, 74}     % untuk blok contoh
\definecolor{ubbred}{RGB}{168, 42, 42}       % untuk blok peringatan
\definecolor{ubbgray}{RGB}{70, 76, 84}       % teks sekunder

% =====================================================================
%  IDENTITAS MATA KULIAH
% =====================================================================
\newcommand{\CourseName}{Pengantar Kriptografi}
\newcommand{\CourseCode}{MPFA014}
\newcommand{\CourseSemester}{Semester Ganjil}
\newcommand{\CourseProgram}{Program Studi Matematika}
\newcommand{\CourseFaculty}{Fakultas Matematika dan Ilmu Pengetahuan Alam}
\newcommand{\CourseInstitution}{Universitas Bale Bandung}
\newcommand{\CourseLecturer}{Cecep Suwanda, S.Si., M.Kom.}

% Label minggu, disetel tiap berkas lewat \setWeekInfo
\newcommand{\WeekLabel}{}
\newcommand{\WeekTopic}{}
\newcommand{\setWeekInfo}[2]{%
    \renewcommand{\WeekLabel}{#1}%
    \renewcommand{\WeekTopic}{#2}%
}

% =====================================================================
%  TEMA: header, footer, warna blok
% =====================================================================
\usetheme{default}
\setbeamertemplate{navigation symbols}{}
\setbeamertemplate{blocks}[rounded][shadow=false]

\setbeamercolor{palette primary}{bg=ubbblue, fg=white}
\setbeamercolor{palette secondary}{bg=ubbgray, fg=white}
\setbeamercolor{palette tertiary}{bg=ubbteal, fg=white}

\setbeamercolor{structure}{fg=ubbblue}
\setbeamercolor{normal text}{fg=black!92}
\setbeamercolor{alerted text}{fg=ubbred}

\setbeamercolor{block title}{bg=ubbblue, fg=white}
\setbeamercolor{block body}{bg=ubbblue!6, fg=black!92}
\setbeamercolor{block title alerted}{bg=ubbred, fg=white}
\setbeamercolor{block body alerted}{bg=ubbred!6, fg=black!92}
\setbeamercolor{block title example}{bg=ubbgreen, fg=white}
\setbeamercolor{block body example}{bg=ubbgreen!6, fg=black!92}

\setbeamercolor{itemize item}{fg=ubbblue}
\setbeamercolor{itemize subitem}{fg=ubbteal}
\setbeamercolor{enumerate item}{fg=ubbblue}

\setbeamerfont{headline}{size=\scriptsize}
\setbeamerfont{footline}{size=\scriptsize}
\setbeamerfont{block title}{size=\normalsize}

% Badan blok rata kiri (bukan rata kanan-kiri): teks slide lebih mudah dibaca
% dan baris terakhir yang pendek tidak lagi memicu "Underfull \hbox".
\addtobeamertemplate{block body begin}{\raggedright}{}
\addtobeamertemplate{block alerted body begin}{\raggedright}{}
\addtobeamertemplate{block example body begin}{\raggedright}{}

% Header: nama mata kuliah dan kode MK
\setbeamertemplate{headline}{%
    \begin{beamercolorbox}[wd=\paperwidth,ht=2.6ex,dp=1.2ex,leftskip=1em,rightskip=1em]{palette primary}%
        \usebeamerfont{headline}\CourseName\ \textperiodcentered\ \CourseCode%
        \hfill \CourseProgram%
    \end{beamercolorbox}%
}

% Footer: minggu dan nomor halaman
\setbeamertemplate{footline}{%
    \begin{beamercolorbox}[wd=\paperwidth,ht=2.4ex,dp=1.2ex,leftskip=1em,rightskip=1em]{palette secondary}%
        \usebeamerfont{footline}\WeekLabel\ \textperiodcentered\ \WeekTopic%
        \hfill \insertframenumber\,/\,\inserttotalframenumber%
    \end{beamercolorbox}%
}

% Halaman judul: judul, subjudul, lalu blok identitas mata kuliah
\setbeamertemplate{title page}{%
    \begin{centering}
        \vskip1.0em
        {\usebeamerfont{title}\usebeamercolor[fg]{title}\inserttitle\par}
        \vskip0.3em
        {\usebeamerfont{subtitle}\usebeamercolor[fg]{subtitle}\insertsubtitle\par}
        \vskip1.2em
        \textcolor{ubbblue}{\rule{0.72\textwidth}{0.8pt}}
        \vskip1.2em
        {\small
            \textbf{\CourseName}\ (\CourseCode)\par
            \vskip0.4em
            \CourseProgram\par
            \CourseFaculty\par
            \CourseInstitution\par
            \vskip0.6em
            \CourseSemester\ \textperiodcentered\ \WeekLabel\par
            \vskip0.8em
            Dosen Pengampu: \CourseLecturer\par
        }
        \vskip1.0em
        \textcolor{ubbblue}{\rule{0.72\textwidth}{0.4pt}}
    \end{centering}
}

% =====================================================================
%  PERINTAH BANTU
% =====================================================================
\newcommand{\R}{\mathbb{R}}
\newcommand{\N}{\mathbb{N}}
\newcommand{\Z}{\mathbb{Z}}
\newcommand{\code}[1]{\texttt{#1}}

% Rujukan gambar dari buku: \bookfig[lebar]{NN}{namaberkas.png}
% Jalur relatif terhadap folder ppt/ (dikompilasi dari sana).
\newcommand{\bookfig}[3][0.80\linewidth]{%
    \includegraphics[width=#1]{../buku_ajar/chapters/#2/figures/#3}%
}

% Keterangan sumber gambar, sepadan dengan \sumber di buku
\newcommand{\figsumber}[1]{%
    \par\vspace{1pt}{\tiny\itshape Sumber: #1\par}%
}

% Penanda blok materi yang sering dipakai di slide
\newcommand{\slidejudul}[1]{{\large\bfseries\textcolor{ubbblue}{#1}}\par\vspace{2pt}}
\newcommand{\kunci}[1]{\textcolor{ubbred}{\textbf{#1}}}

%      \title{...} \subtitle{...} \setWeekInfo{Minggu n}{Topik}
%
%  Berkas minggu hanya berisi \setWeekInfo dan lingkungan frame.
%  Kompilasi dijalankan DARI DALAM folder ppt/ (lihat compile-ppt.bat),
%  karena gambar buku dirujuk dengan jalur relatif ../buku_ajar/.
% =====================================================================

% ---------------------------------------------------------------- bahasa
\usepackage[utf8]{inputenc}
\usepackage[T1]{fontenc}
\usepackage[indonesian]{babel}
\usepackage{lmodern}
\usepackage{microtype}

% ------------------------------------------------------------- matematika
\usepackage{amsmath,amssymb,mathtools}

% ----------------------------------------------------------------- tabel
\usepackage{booktabs}
\usepackage{array}
\usepackage{tabularx}

% --------------------------------------------------------------- grafik
\usepackage{tikz}
\usetikzlibrary{arrows.meta,positioning,shapes.geometric,calc,fit,backgrounds}

% ---------------------------------------------------------- kode program
\usepackage{listings}
\lstdefinestyle{slidecode}{
    basicstyle=\ttfamily\scriptsize,
    keywordstyle=\color{ubbblue}\bfseries,
    commentstyle=\color{ubbgreen}\itshape,
    stringstyle=\color{ubbred},
    numbers=left,
    numberstyle=\tiny\color{gray},
    numbersep=6pt,
    frame=single,
    rulecolor=\color{ubbblue!40},
    backgroundcolor=\color{ubbblue!3},
    breaklines=true,
    showstringspaces=false,
    tabsize=4,
    columns=fullflexible,
    keepspaces=true,
}
\lstset{style=slidecode}

% =====================================================================
%  WARNA INSTITUSI
% =====================================================================
\definecolor{ubbblue}{RGB}{12, 62, 120}      % warna utama
\definecolor{ubbteal}{RGB}{0, 118, 110}      % aksen kedua
\definecolor{ubbgreen}{RGB}{22, 122, 74}     % untuk blok contoh
\definecolor{ubbred}{RGB}{168, 42, 42}       % untuk blok peringatan
\definecolor{ubbgray}{RGB}{70, 76, 84}       % teks sekunder

% =====================================================================
%  IDENTITAS MATA KULIAH
% =====================================================================
\newcommand{\CourseName}{Pengantar Kriptografi}
\newcommand{\CourseCode}{MPFA014}
\newcommand{\CourseSemester}{Semester Ganjil}
\newcommand{\CourseProgram}{Program Studi Matematika}
\newcommand{\CourseFaculty}{Fakultas Matematika dan Ilmu Pengetahuan Alam}
\newcommand{\CourseInstitution}{Universitas Bale Bandung}
\newcommand{\CourseLecturer}{Cecep Suwanda, S.Si., M.Kom.}

% Label minggu, disetel tiap berkas lewat \setWeekInfo
\newcommand{\WeekLabel}{}
\newcommand{\WeekTopic}{}
\newcommand{\setWeekInfo}[2]{%
    \renewcommand{\WeekLabel}{#1}%
    \renewcommand{\WeekTopic}{#2}%
}

% =====================================================================
%  TEMA: header, footer, warna blok
% =====================================================================
\usetheme{default}
\setbeamertemplate{navigation symbols}{}
\setbeamertemplate{blocks}[rounded][shadow=false]

\setbeamercolor{palette primary}{bg=ubbblue, fg=white}
\setbeamercolor{palette secondary}{bg=ubbgray, fg=white}
\setbeamercolor{palette tertiary}{bg=ubbteal, fg=white}

\setbeamercolor{structure}{fg=ubbblue}
\setbeamercolor{normal text}{fg=black!92}
\setbeamercolor{alerted text}{fg=ubbred}

\setbeamercolor{block title}{bg=ubbblue, fg=white}
\setbeamercolor{block body}{bg=ubbblue!6, fg=black!92}
\setbeamercolor{block title alerted}{bg=ubbred, fg=white}
\setbeamercolor{block body alerted}{bg=ubbred!6, fg=black!92}
\setbeamercolor{block title example}{bg=ubbgreen, fg=white}
\setbeamercolor{block body example}{bg=ubbgreen!6, fg=black!92}

\setbeamercolor{itemize item}{fg=ubbblue}
\setbeamercolor{itemize subitem}{fg=ubbteal}
\setbeamercolor{enumerate item}{fg=ubbblue}

\setbeamerfont{headline}{size=\scriptsize}
\setbeamerfont{footline}{size=\scriptsize}
\setbeamerfont{block title}{size=\normalsize}

% Badan blok rata kiri (bukan rata kanan-kiri): teks slide lebih mudah dibaca
% dan baris terakhir yang pendek tidak lagi memicu "Underfull \hbox".
\addtobeamertemplate{block body begin}{\raggedright}{}
\addtobeamertemplate{block alerted body begin}{\raggedright}{}
\addtobeamertemplate{block example body begin}{\raggedright}{}

% Header: nama mata kuliah dan kode MK
\setbeamertemplate{headline}{%
    \begin{beamercolorbox}[wd=\paperwidth,ht=2.6ex,dp=1.2ex,leftskip=1em,rightskip=1em]{palette primary}%
        \usebeamerfont{headline}\CourseName\ \textperiodcentered\ \CourseCode%
        \hfill \CourseProgram%
    \end{beamercolorbox}%
}

% Footer: minggu dan nomor halaman
\setbeamertemplate{footline}{%
    \begin{beamercolorbox}[wd=\paperwidth,ht=2.4ex,dp=1.2ex,leftskip=1em,rightskip=1em]{palette secondary}%
        \usebeamerfont{footline}\WeekLabel\ \textperiodcentered\ \WeekTopic%
        \hfill \insertframenumber\,/\,\inserttotalframenumber%
    \end{beamercolorbox}%
}

% Halaman judul: judul, subjudul, lalu blok identitas mata kuliah
\setbeamertemplate{title page}{%
    \begin{centering}
        \vskip1.0em
        {\usebeamerfont{title}\usebeamercolor[fg]{title}\inserttitle\par}
        \vskip0.3em
        {\usebeamerfont{subtitle}\usebeamercolor[fg]{subtitle}\insertsubtitle\par}
        \vskip1.2em
        \textcolor{ubbblue}{\rule{0.72\textwidth}{0.8pt}}
        \vskip1.2em
        {\small
            \textbf{\CourseName}\ (\CourseCode)\par
            \vskip0.4em
            \CourseProgram\par
            \CourseFaculty\par
            \CourseInstitution\par
            \vskip0.6em
            \CourseSemester\ \textperiodcentered\ \WeekLabel\par
            \vskip0.8em
            Dosen Pengampu: \CourseLecturer\par
        }
        \vskip1.0em
        \textcolor{ubbblue}{\rule{0.72\textwidth}{0.4pt}}
    \end{centering}
}

% =====================================================================
%  PERINTAH BANTU
% =====================================================================
\newcommand{\R}{\mathbb{R}}
\newcommand{\N}{\mathbb{N}}
\newcommand{\Z}{\mathbb{Z}}
\newcommand{\code}[1]{\texttt{#1}}

% Rujukan gambar dari buku: \bookfig[lebar]{NN}{namaberkas.png}
% Jalur relatif terhadap folder ppt/ (dikompilasi dari sana).
\newcommand{\bookfig}[3][0.80\linewidth]{%
    \includegraphics[width=#1]{../buku_ajar/chapters/#2/figures/#3}%
}

% Keterangan sumber gambar, sepadan dengan \sumber di buku
\newcommand{\figsumber}[1]{%
    \par\vspace{1pt}{\tiny\itshape Sumber: #1\par}%
}

% Penanda blok materi yang sering dipakai di slide
\newcommand{\slidejudul}[1]{{\large\bfseries\textcolor{ubbblue}{#1}}\par\vspace{2pt}}
\newcommand{\kunci}[1]{\textcolor{ubbred}{\textbf{#1}}}

%      \title{...} \subtitle{...} \setWeekInfo{Minggu n}{Topik}
%
%  Berkas minggu hanya berisi \setWeekInfo dan lingkungan frame.
%  Kompilasi dijalankan DARI DALAM folder ppt/ (lihat compile-ppt.bat),
%  karena gambar buku dirujuk dengan jalur relatif ../buku_ajar/.
% =====================================================================

% ---------------------------------------------------------------- bahasa
\usepackage[utf8]{inputenc}
\usepackage[T1]{fontenc}
\usepackage[indonesian]{babel}
\usepackage{lmodern}
\usepackage{microtype}

% ------------------------------------------------------------- matematika
\usepackage{amsmath,amssymb,mathtools}

% ----------------------------------------------------------------- tabel
\usepackage{booktabs}
\usepackage{array}
\usepackage{tabularx}

% --------------------------------------------------------------- grafik
\usepackage{tikz}
\usetikzlibrary{arrows.meta,positioning,shapes.geometric,calc,fit,backgrounds}

% ---------------------------------------------------------- kode program
\usepackage{listings}
\lstdefinestyle{slidecode}{
    basicstyle=\ttfamily\scriptsize,
    keywordstyle=\color{ubbblue}\bfseries,
    commentstyle=\color{ubbgreen}\itshape,
    stringstyle=\color{ubbred},
    numbers=left,
    numberstyle=\tiny\color{gray},
    numbersep=6pt,
    frame=single,
    rulecolor=\color{ubbblue!40},
    backgroundcolor=\color{ubbblue!3},
    breaklines=true,
    showstringspaces=false,
    tabsize=4,
    columns=fullflexible,
    keepspaces=true,
}
\lstset{style=slidecode}

% =====================================================================
%  WARNA INSTITUSI
% =====================================================================
\definecolor{ubbblue}{RGB}{12, 62, 120}      % warna utama
\definecolor{ubbteal}{RGB}{0, 118, 110}      % aksen kedua
\definecolor{ubbgreen}{RGB}{22, 122, 74}     % untuk blok contoh
\definecolor{ubbred}{RGB}{168, 42, 42}       % untuk blok peringatan
\definecolor{ubbgray}{RGB}{70, 76, 84}       % teks sekunder

% =====================================================================
%  IDENTITAS MATA KULIAH
% =====================================================================
\newcommand{\CourseName}{Pengantar Kriptografi}
\newcommand{\CourseCode}{MPFA014}
\newcommand{\CourseSemester}{Semester Ganjil}
\newcommand{\CourseProgram}{Program Studi Matematika}
\newcommand{\CourseFaculty}{Fakultas Matematika dan Ilmu Pengetahuan Alam}
\newcommand{\CourseInstitution}{Universitas Bale Bandung}
\newcommand{\CourseLecturer}{Cecep Suwanda, S.Si., M.Kom.}

% Label minggu, disetel tiap berkas lewat \setWeekInfo
\newcommand{\WeekLabel}{}
\newcommand{\WeekTopic}{}
\newcommand{\setWeekInfo}[2]{%
    \renewcommand{\WeekLabel}{#1}%
    \renewcommand{\WeekTopic}{#2}%
}

% =====================================================================
%  TEMA: header, footer, warna blok
% =====================================================================
\usetheme{default}
\setbeamertemplate{navigation symbols}{}
\setbeamertemplate{blocks}[rounded][shadow=false]

\setbeamercolor{palette primary}{bg=ubbblue, fg=white}
\setbeamercolor{palette secondary}{bg=ubbgray, fg=white}
\setbeamercolor{palette tertiary}{bg=ubbteal, fg=white}

\setbeamercolor{structure}{fg=ubbblue}
\setbeamercolor{normal text}{fg=black!92}
\setbeamercolor{alerted text}{fg=ubbred}

\setbeamercolor{block title}{bg=ubbblue, fg=white}
\setbeamercolor{block body}{bg=ubbblue!6, fg=black!92}
\setbeamercolor{block title alerted}{bg=ubbred, fg=white}
\setbeamercolor{block body alerted}{bg=ubbred!6, fg=black!92}
\setbeamercolor{block title example}{bg=ubbgreen, fg=white}
\setbeamercolor{block body example}{bg=ubbgreen!6, fg=black!92}

\setbeamercolor{itemize item}{fg=ubbblue}
\setbeamercolor{itemize subitem}{fg=ubbteal}
\setbeamercolor{enumerate item}{fg=ubbblue}

\setbeamerfont{headline}{size=\scriptsize}
\setbeamerfont{footline}{size=\scriptsize}
\setbeamerfont{block title}{size=\normalsize}

% Badan blok rata kiri (bukan rata kanan-kiri): teks slide lebih mudah dibaca
% dan baris terakhir yang pendek tidak lagi memicu "Underfull \hbox".
\addtobeamertemplate{block body begin}{\raggedright}{}
\addtobeamertemplate{block alerted body begin}{\raggedright}{}
\addtobeamertemplate{block example body begin}{\raggedright}{}

% Header: nama mata kuliah dan kode MK
\setbeamertemplate{headline}{%
    \begin{beamercolorbox}[wd=\paperwidth,ht=2.6ex,dp=1.2ex,leftskip=1em,rightskip=1em]{palette primary}%
        \usebeamerfont{headline}\CourseName\ \textperiodcentered\ \CourseCode%
        \hfill \CourseProgram%
    \end{beamercolorbox}%
}

% Footer: minggu dan nomor halaman
\setbeamertemplate{footline}{%
    \begin{beamercolorbox}[wd=\paperwidth,ht=2.4ex,dp=1.2ex,leftskip=1em,rightskip=1em]{palette secondary}%
        \usebeamerfont{footline}\WeekLabel\ \textperiodcentered\ \WeekTopic%
        \hfill \insertframenumber\,/\,\inserttotalframenumber%
    \end{beamercolorbox}%
}

% Halaman judul: judul, subjudul, lalu blok identitas mata kuliah
\setbeamertemplate{title page}{%
    \begin{centering}
        \vskip1.0em
        {\usebeamerfont{title}\usebeamercolor[fg]{title}\inserttitle\par}
        \vskip0.3em
        {\usebeamerfont{subtitle}\usebeamercolor[fg]{subtitle}\insertsubtitle\par}
        \vskip1.2em
        \textcolor{ubbblue}{\rule{0.72\textwidth}{0.8pt}}
        \vskip1.2em
        {\small
            \textbf{\CourseName}\ (\CourseCode)\par
            \vskip0.4em
            \CourseProgram\par
            \CourseFaculty\par
            \CourseInstitution\par
            \vskip0.6em
            \CourseSemester\ \textperiodcentered\ \WeekLabel\par
            \vskip0.8em
            Dosen Pengampu: \CourseLecturer\par
        }
        \vskip1.0em
        \textcolor{ubbblue}{\rule{0.72\textwidth}{0.4pt}}
    \end{centering}
}

% =====================================================================
%  PERINTAH BANTU
% =====================================================================
\newcommand{\R}{\mathbb{R}}
\newcommand{\N}{\mathbb{N}}
\newcommand{\Z}{\mathbb{Z}}
\newcommand{\code}[1]{\texttt{#1}}

% Rujukan gambar dari buku: \bookfig[lebar]{NN}{namaberkas.png}
% Jalur relatif terhadap folder ppt/ (dikompilasi dari sana).
\newcommand{\bookfig}[3][0.80\linewidth]{%
    \includegraphics[width=#1]{../buku_ajar/chapters/#2/figures/#3}%
}

% Keterangan sumber gambar, sepadan dengan \sumber di buku
\newcommand{\figsumber}[1]{%
    \par\vspace{1pt}{\tiny\itshape Sumber: #1\par}%
}

% Penanda blok materi yang sering dipakai di slide
\newcommand{\slidejudul}[1]{{\large\bfseries\textcolor{ubbblue}{#1}}\par\vspace{2pt}}
\newcommand{\kunci}[1]{\textcolor{ubbred}{\textbf{#1}}}


\title{Layanan Keamanan dan Serangan}
\subtitle{Membaca serangan dari informasi, tujuan, dan keterlibatan penyerang}
\setWeekInfo{Minggu 2}{Layanan keamanan \& serangan}

\begin{document}

% ----------------------------------------------------------- halaman judul
\begin{frame}[plain]
    \titlepage
\end{frame}

% ------------------------------------------------------------------ peta
\begin{frame}{Peta Materi Pertemuan Ini}
    \begin{columns}[T,onlytextwidth]
        \begin{column}{0.52\textwidth}
            \slidejudul{Lima bagian}
            \begin{enumerate}
                \item Konsep dasar serangan
                \item Jenis serangan berdasarkan informasi
                \item Tujuan serangan
                \item Serangan pasif dan serangan aktif
                \item Keamanan komputasi dan serangan saluran samping
            \end{enumerate}
        \end{column}
        \begin{column}{0.44\textwidth}
            \begin{block}{Referensi}
                \footnotesize
                Bab 2 buku ajar; Munir, \textit{IF4020 Kriptografi},
                ITB, 2026 (modul 06).
            \end{block}
            \begin{exampleblock}{Tiga pertanyaan pemandu}
                \footnotesize
                Apa yang \emph{dimiliki} penyerang? Apa yang ingin ia
                \emph{peroleh}? Apakah ia \emph{mengubah} komunikasi?
            \end{exampleblock}
        \end{column}
    \end{columns}
\end{frame}

% ---------------------------------------------------- capaian pembelajaran
\begin{frame}{Capaian Pembelajaran}
    \footnotesize
    \begin{block}{Sub-CPMK yang dibebankan}
        \begin{itemize}
            \item \textbf{Sub-CPMK 2.1} --- Menjelaskan definisi, terminologi,
                  sejarah, dan empat layanan kriptografi (kerahasiaan, integritas,
                  autentikasi, dan non-repudiation).
            \item \textbf{Sub-CPMK 2.4} --- Mengidentifikasi jenis-jenis serangan
                  terhadap kriptografi berdasarkan data yang tersedia, sifat
                  serangan pasif atau aktif, dan \textit{side-channel attack}.
        \end{itemize}
    \end{block}
    \vspace{0.2cm}
    \begin{exampleblock}{Setelah pertemuan ini, mahasiswa mampu}
        \begin{itemize}
            \item menjelaskan definisi serangan dan Prinsip Kerckhoffs;
            \item mengidentifikasi lima jenis serangan berdasarkan informasi
                  yang tersedia;
            \item membedakan serangan pasif dan aktif beserta contohnya;
            \item mengidentifikasi empat jenis serangan saluran samping.
        \end{itemize}
    \end{exampleblock}
    \vspace{0.1cm}
    {\scriptsize\textcolor{ubbgray}{Cakupan 2.1 pada pertemuan ini: definisi
    dan terminologi serangan; empat layanan telah dibahas pada Minggu 1.}}
\end{frame}

% =====================================================================
\section{Konsep Dasar Serangan}
% =====================================================================

\begin{frame}{Apa yang Disebut Serangan?}
    \begin{block}{Definisi}
        Serangan (\textit{attack}) adalah setiap usaha pihak lawan untuk
        menemukan kunci atau menemukan plainteks dari cipherteksnya.
    \end{block}
    \vspace{0.25cm}
    \begin{columns}[T,onlytextwidth]
        \begin{column}{0.48\textwidth}
            \slidejudul{Pihak lawan}
            \begin{itemize}
                \item Disebut musuh, lawan, penyadap, atau kriptanalis.
                \item Bertujuan memperoleh plainteks atau kunci dekripsi.
            \end{itemize}
        \end{column}
        \begin{column}{0.48\textwidth}
            \slidejudul{Mengapa kunci diburu}
            \begin{itemize}
                \item Satu kunci membuka potongan cipherteks yang lain.
                \item Dengan kunci itu, cipherteks berikutnya dapat
                      didekripsi.
            \end{itemize}
        \end{column}
    \end{columns}
    \vspace{0.2cm}
    \begin{alertblock}{Tujuan kriptografi}
        Menjaga kerahasiaan pesan atau kunci dari pihak ketiga.
    \end{alertblock}
\end{frame}

\begin{frame}{Prinsip Kerckhoffs}
    \begin{alertblock}{Prinsip Kerckhoffs}
        Semua algoritma kriptografi harus publik; hanya kuncinya yang rahasia.
    \end{alertblock}
    \vspace{0.3cm}
    \begin{columns}[T,onlytextwidth]
        \begin{column}{0.46\textwidth}
            \begin{block}{Yang publik}
                Algoritma dan cara kerjanya.
            \end{block}
        \end{column}
        \begin{column}{0.46\textwidth}
            \begin{block}{Yang rahasia}
                Kunci yang dipakai untuk enkripsi dan dekripsi.
            \end{block}
        \end{column}
    \end{columns}
    \vspace{0.3cm}
    \begin{exampleblock}{Konsekuensi}
        Penyerang dianggap mengetahui algoritma yang dipakai. Keamanan sistem
        terletak pada \kunci{kunci}, bukan pada kerahasiaan algoritmanya.
    \end{exampleblock}
\end{frame}

\begin{frame}{Tiga Cara Mengelompokkan Serangan}
    Serangan dapat dibaca dari tiga sudut yang berbeda.
    \vspace{0.35cm}
    \begin{enumerate}
        \item Berdasarkan \textbf{informasi yang tersedia} bagi penyerang.
              \hfill\textcolor{ubbblue}{\textbf{apa yang dimilikinya?}}
        \item Berdasarkan \textbf{tujuan atau hasil} yang ingin dicapai.
              \hfill\textcolor{ubbblue}{\textbf{apa yang ingin diperoleh?}}
        \item Berdasarkan \textbf{keterlibatan penyerang} dalam komunikasi.
              \hfill\textcolor{ubbblue}{\textbf{apakah ia mengubah data?}}
    \end{enumerate}
    \vspace{0.35cm}
    \begin{exampleblock}{Cara membaca}
        Ketiganya adalah lensa yang berbeda untuk serangan yang sama. Satu
        serangan dapat dijelaskan dari lebih dari satu sudut.
    \end{exampleblock}
\end{frame}

% =====================================================================
\section{Serangan Berdasarkan Informasi}
% =====================================================================

\begin{frame}{Lima Jenis Berdasarkan Informasi}
    \scriptsize
    Semakin banyak informasi di tangan penyerang, semakin besar peluang
    serangan berhasil.
    \vspace{0.1cm}
    \begin{tabularx}{\linewidth}{@{}>{\raggedright\arraybackslash}p{2.45cm}
        >{\raggedright\arraybackslash}p{4.55cm}
        >{\raggedright\arraybackslash}X@{}}
        \toprule
        \textbf{Jenis} & \textbf{Informasi yang dimiliki} & \textbf{Ilustrasi pada Caesar} \\
        \midrule
        \textit{Ciphertext-only} & Hanya cipherteks $C_i=E_k(P_i)$. &
        \code{KHOORZRUOG}; tebak $P$ atau $k$. \\
        \addlinespace
        \textit{Known-plaintext} & Pasangan $P_i$ dan $C_i=E_k(P_i)$. &
        Duga ada \code{HELLO}; dari H$\to$K, $k=3$. \\
        \addlinespace
        \textit{Chosen-plaintext} & Memilih $P_i$; punya akses ke sistem
        enkripsi. & Kirim \code{AAAAA}; keluar \code{FFFFF}, $k=5$. \\
        \addlinespace
        \textit{Adaptive chosen-plaintext} & $P$ berikutnya menyesuaikan $C$
        sebelumnya. & \code{A}$\to$\code{F}, lalu \code{B}$\to$\code{G};
        $k=5$. \\
        \addlinespace
        \textit{Chosen-ciphertext} & Memilih $C_i$; memakai \textit{oracle}
        dekripsi. & \code{A} didekripsi menjadi \code{X}; $k=3$. \\
        \bottomrule
    \end{tabularx}
\end{frame}

\begin{frame}{\textit{Ciphertext-Only Attack}}
    \begin{columns}[T,onlytextwidth]
        \begin{column}{0.54\textwidth}
            \slidejudul{Hanya cipherteks}
            \begin{itemize}
                \item Informasi paling sedikit, sehingga serangan paling sulit.
                \item Penyerang memegang $C_j=E_k(P_j)$ untuk $j=1,\dots,i$.
                \item Ia harus mendeduksi $P_1,\dots,P_i$ atau kunci $k$.
            \end{itemize}
            \begin{exampleblock}{Teknik yang dipakai}
                \footnotesize
                Pencarian kunci menyeluruh, terkaan, dan analisis frekuensi.
            \end{exampleblock}
        \end{column}
        \begin{column}{0.44\textwidth}
            \centering
            \bookfig[0.95\linewidth]{02}{serangan-ciphertext-only.png}
            \figsumber{Munir, \textit{IF4020 Kriptografi}, ITB.}
        \end{column}
    \end{columns}
    \vspace{0.05cm}
    \begin{block}{Ilustrasi dari bab}
        \footnotesize
        Caesar \code{KHOORZRUOG}: ada 25 kemungkinan pergeseran; percobaan
        $k=3$ menghasilkan \code{HELLOWORLD}.
    \end{block}
\end{frame}

\begin{frame}{\textit{Known-Plaintext Attack}}
    \begin{columns}[T,onlytextwidth]
        \begin{column}{0.55\textwidth}
            \slidejudul{Pasangan yang diketahui}
            \begin{itemize}
                \item Penyerang memiliki pasangan $P_j$ dan $C_j=E_k(P_j)$
                      untuk $j=1,\dots,i$.
                \item Yang dicari adalah kunci $k$.
                \item Dengan $k$, $P_{i+1}$ diperoleh dari
                      $C_{i+1}=E_k(P_{i+1})$.
            \end{itemize}
        \end{column}
        \begin{column}{0.43\textwidth}
            \begin{exampleblock}{Mengapa sering terjadi}
                \footnotesize
                \begin{itemize}
                    \item Surel memuat \textit{From} dan \textit{To}.
                    \item Surat resmi memuat ``Dengan hormat'' atau
                          ``wassalam''.
                    \item Berkas sumber memuat \code{\#include} atau
                          \code{program}.
                \end{itemize}
            \end{exampleblock}
        \end{column}
    \end{columns}
    \vspace{0.1cm}
    \begin{alertblock}{Contoh pada Caesar}
        \footnotesize
        Cipherteks \code{KHOOR ZRUOG}; penyerang menduga ada \code{HELLO}.
        Dari pasangan H$\to$K ia menyimpulkan $k=3$, lalu membuka seluruh
        cipherteks.
    \end{alertblock}
\end{frame}

\begin{frame}{\textit{Chosen-Plaintext Attack}}
    \begin{columns}[T,onlytextwidth]
        \begin{column}{0.55\textwidth}
            \slidejudul{Memilih plainteks}
            \begin{itemize}
                \item Penyerang dapat memilih plainteks yang akan
                      dienkripsikan.
                \item Ia dianggap memiliki akses ke sistem enkripsi.
                \item Diperoleh pasangan $P_j$ dan $C_j=E_k(P_j)$ untuk
                      $j=1,\dots,i$.
                \item Target akhirnya tetap kunci $k$.
                \item Contoh Caesar: pilih \code{AAAAA}, keluar
                      \code{FFFFF}; A$\to$F berarti $k=5$.
            \end{itemize}
        \end{column}
        \begin{column}{0.43\textwidth}
            \centering
            \bookfig[0.98\linewidth]{02}{serangan-chosen-plaintext.png}
            \figsumber{Munir, \textit{IF4020 Kriptografi}, ITB.}
        \end{column}
    \end{columns}
    \vspace{0.05cm}
    \begin{block}{Kasus pada mesin ATM}
        \footnotesize
        Kriptanalis 1 mengubah PIN; kriptanalis 2 menyadap cipherteks bank.
    \end{block}
\end{frame}

\begin{frame}{Contoh Terhitung: \textit{Chosen-Plaintext} pada Vigen\`ere}
    \footnotesize
    Pada Vigen\`ere berlaku $C_i=(P_i+K_i)\bmod 26$.
    \vspace{0.1cm}
    \begin{enumerate}
        \item Kunci \code{SCRAM}: S$\to$18, C$\to$2, R$\to$17, A$\to$0,
              M$\to$12, sehingga $K=(18,2,17,0,12)$.
        \item Penyerang memilih $P=\code{AAAAA}$, yaitu
              $(0,0,0,0,0)$.
        \item Hitung $C_i=(0+K_i)\bmod 26=K_i$.
        \item Hasil $C=(18,2,17,0,12)=\code{SCRAM}$.
    \end{enumerate}
    \vspace{0.1cm}
    \begin{alertblock}{Hasil}
        Cipherteks yang keluar sama dengan kuncinya. Satu \textit{query} sudah
        cukup untuk memulihkan seluruh kunci.
    \end{alertblock}
\end{frame}

\begin{frame}{\textit{Adaptive Chosen-Plaintext Attack}}
    \slidejudul{Variasi yang menyesuaikan diri}
    \begin{itemize}
        \item Masih termasuk \textit{chosen-plaintext attack}.
        \item Penyerang memilih plainteks berikutnya berdasarkan cipherteks
              dari \textit{query} sebelumnya.
        \item Ia belajar secara bertahap, bukan menyusun seluruh
              \textit{query} di muka.
    \end{itemize}
    \vspace{0.15cm}
    \begin{columns}[T,onlytextwidth]
        \begin{column}{0.47\textwidth}
            \begin{block}{Contoh pada Caesar}
                \footnotesize
                Pilih \code{A}, keluar \code{F}; pergeseran mungkin 5.
                Pilih \code{B}, keluar \code{G}; polanya konsisten, sehingga
                $k=5$ dipastikan.
            \end{block}
        \end{column}
        \begin{column}{0.49\textwidth}
            \begin{exampleblock}{Contoh pada Vigen\`ere}
                \footnotesize
                Masukkan \code{AAAA} sehingga cipherteks sama dengan kunci.
                Setelah kunci diketahui, masukkan \code{BBBB} untuk
                memverifikasi hasilnya.
            \end{exampleblock}
        \end{column}
    \end{columns}
\end{frame}

\begin{frame}{\textit{Chosen-Ciphertext Attack}}
    \begin{columns}[T,onlytextwidth]
        \begin{column}{0.56\textwidth}
            \slidejudul{Oracle dekripsi}
            \begin{itemize}
                \item Penyerang memilih cipherteks tertentu.
                \item Ia memakai \textit{oracle}, yaitu sistem yang dapat
                      mendekripsi.
                \item Diperoleh pasangan $C_j$ dan $P_j=D_k(C_j)$ untuk
                      $j=1,\dots,i$.
                \item Hasil dekripsi dipakai untuk menemukan kunci.
            \end{itemize}
        \end{column}
        \begin{column}{0.40\textwidth}
            \begin{exampleblock}{Contoh pada Caesar}
                \footnotesize
                Pilih $C=\code{A}$; dekripsi menghasilkan $P=\code{X}$.
                Karena A adalah X yang digeser, $k=3$.
            \end{exampleblock}
        \end{column}
    \end{columns}
    \vspace{0.15cm}
    \begin{alertblock}{Inti}
        Penyerang tidak menebak plainteks; ia mengirim cipherteks dan membaca
        jawaban sistem.
    \end{alertblock}
\end{frame}

% =====================================================================
\section{Tujuan Serangan}
% =====================================================================

\begin{frame}{Lima Tujuan Serangan}
    \footnotesize
    Pertanyaannya berubah: bukan ``apa yang dimiliki penyerang?'', melainkan
    ``apa yang ingin ia capai?''
    \vspace{0.1cm}
    \begin{tabularx}{\linewidth}{@{}>{\raggedright\arraybackslash}p{3.1cm}
        >{\raggedright\arraybackslash}X@{}}
        \toprule
        \textbf{Tujuan} & \textbf{Yang ingin dicapai} \\
        \midrule
        \textit{Key-recovery} & Memperoleh kunci $K$ pada $C=E_K(P)$. \\
        \addlinespace
        \textit{Plaintext-recovery} & Memperoleh plainteks, meskipun kunci
        tidak ditemukan. \\
        \addlinespace
        \textit{Distinguishing} & Menentukan apakah data adalah keluaran
        algoritma kriptografi atau data acak. \\
        \addlinespace
        \textit{Forgery} & Membuat pesan palsu yang diterima sistem sebagai
        pesan sah. \\
        \addlinespace
        \textit{Replay} & Mengirim ulang pesan sah untuk memicu tindakan yang
        sama. \\
        \bottomrule
    \end{tabularx}
\end{frame}

\begin{frame}{Ukuran Keberhasilan Tiga Tujuan}
    \begin{columns}[T,onlytextwidth]
        \begin{column}{0.48\textwidth}
            \begin{block}{\textit{Key-recovery}}
                \footnotesize
                Kunci 16 karakter dicoba satu per satu:
                $K_1\to P_1$, $K_2\to P_2$, $K_3\to P_3$, \dots sampai
                ditemukan plainteks bermakna.
            \end{block}
            \begin{block}{\textit{Plaintext-recovery}}
                \footnotesize
                Cipherteks $C=\text{XQZ}\dots$; strukturnya menunjukkan pesan
                mungkin berbunyi \text{MEET AT THE STATION}.
            \end{block}
        \end{column}
        \begin{column}{0.48\textwidth}
            \begin{block}{\textit{Distinguishing}}
                \footnotesize
                Keluaran $10101010\dots$ jelas berpola, sehingga mudah
                dibedakan dari data acak.
            \end{block}
            \begin{alertblock}{\textit{Forgery}}
                \footnotesize
                Membuat pesan palsu yang diterima sistem sebagai pesan sah.
            \end{alertblock}
        \end{column}
    \end{columns}
\end{frame}

\begin{frame}{\textit{Replay Attack} dan Penangkalannya}
    \begin{columns}[T,onlytextwidth]
        \begin{column}{0.55\textwidth}
            \slidejudul{Tanpa memecahkan enkripsi}
            \begin{itemize}
                \item Penyerang menangkap pesan sah lalu mengirimnya ulang.
                \item Setiap pengiriman ulang memicu tindakan yang sama.
                \item Penangkal: \textit{nonce}, cap waktu
                      (\textit{timestamp}), atau nomor urut
                      (\textit{sequence number}).
            \end{itemize}
            \begin{exampleblock}{Nomor urut}
                \footnotesize
                Sistem menyimpan nomor terakhir 103. Perintah 101 ditolak
                karena $101\not>103$; perintah 104 diterima dan menjadi nomor
                terakhir.
            \end{exampleblock}
        \end{column}
        \begin{column}{0.43\textwidth}
            \centering
            \bookfig[0.98\linewidth]{02}{serangan-replay.png}
            \figsumber{Munir, \textit{IF4020 Kriptografi}, ITB.}
        \end{column}
    \end{columns}
\end{frame}

% =====================================================================
\section{Serangan Pasif dan Serangan Aktif}
% =====================================================================

\begin{frame}{Serangan Pasif}
    \begin{columns}[T,onlytextwidth]
        \begin{column}{0.55\textwidth}
            \slidejudul{Tidak terlibat, hanya menyadap}
            \begin{itemize}
                \item Penyerang tidak mengubah komunikasi.
                \item Ia mengumpulkan data sebanyak mungkin.
                \item Enkripsi mencegah penyadap memahami isi pesan.
            \end{itemize}
            \begin{block}{Dua bentuk}
                \begin{itemize}
                    \item \textbf{Intersepsi}: menyadap isi pesan.
                    \item \textbf{Analisis lalu lintas}: mengamati pola tanpa
                          membaca isinya.
                \end{itemize}
            \end{block}
        \end{column}
        \begin{column}{0.43\textwidth}
            \centering
            \bookfig[0.98\linewidth]{02}{serangan-pasif-intersepsi.png}
            \figsumber{Munir, \textit{IF4020 Kriptografi}, ITB.}
        \end{column}
    \end{columns}
\end{frame}

\begin{frame}{Analisis Lalu Lintas}
    \begin{columns}[T,onlytextwidth]
        \begin{column}{0.53\textwidth}
            \slidejudul{Isi tersembunyi, pola tetap terlihat}
            \begin{itemize}
                \item Penyerang mengamati pola lalu lintas.
                \item Ia tidak membaca isi pesan.
                \item Yang menjadi sasaran adalah pola komunikasi, bukan
                      maknanya.
            \end{itemize}
        \end{column}
        \begin{column}{0.45\textwidth}
            \centering
            \bookfig[0.98\linewidth]{02}{serangan-analisis-lalu-lintas.png}
            \figsumber{Munir, \textit{IF4020 Kriptografi}, ITB.}
        \end{column}
    \end{columns}
    \vspace{0.05cm}
    \begin{block}{Poin}
        Pola lalu lintas saja sudah memberi informasi, meskipun isi pesannya
        tidak terbaca.
    \end{block}
\end{frame}

\begin{frame}{Ketika Data Tidak Dienkripsi: Wireshark}
    \begin{columns}[T,onlytextwidth]
        \begin{column}{0.52\textwidth}
            \slidejudul{Isi terbaca langsung}
            \begin{itemize}
                \item Wireshark memantau lalu lintas jaringan.
                \item Bila data tidak dienkripsi, penyadap tidak perlu
                      menganalisis apa pun.
                \item Nama pengguna dan kata sandi tampak pada tangkapan
                      paket.
            \end{itemize}
        \end{column}
        \begin{column}{0.46\textwidth}
            \centering
            \bookfig[1.0\linewidth]{02}{wireshark-plaintext.png}
            \figsumber{Munir, \textit{IF4020 Kriptografi}, ITB.}
        \end{column}
    \end{columns}
\end{frame}

\begin{frame}{Tiga Metode Penyadapan Fisik}
    \begin{columns}[T,onlytextwidth]
        \begin{column}{0.56\textwidth}
            \footnotesize
            Selain penyadapan perangkat lunak, dikenal tiga metode yang
            bekerja pada tingkat fisik.
            \begin{enumerate}
                \item \textbf{Penyadapan kabel} (\textit{wiretapping}):
                      menempelkan alat pada saluran komunikasi. Sistem
                      telepon kabel rentan karena rancangannya sederhana.
                \item \textbf{Penyadapan elektromagnetik}: menangkap radiasi
                      elektromagnetik yang dipancarkan perangkat elektronik.
                \item \textbf{Penyadapan akustik}: memanfaatkan getaran atau
                      suara, termasuk \textit{vibrometry} nirkabel.
            \end{enumerate}
        \end{column}
        \begin{column}{0.42\textwidth}
            \centering
            \bookfig[1.0\linewidth]{02}{penyadapan-elektromagnetik.png}
            \figsumber{Munir, \textit{IF4020 Kriptografi}, ITB.}
        \end{column}
    \end{columns}
\end{frame}

\begin{frame}{Penyadapan Akustik}
    \begin{columns}[T,onlytextwidth]
        \begin{column}{0.50\textwidth}
            \slidejudul{Getaran dan suara}
            \begin{itemize}
                \item Penyadapan akustik memanfaatkan getaran atau suara.
                \item Salah satu ragamnya adalah \textit{vibrometry}
                      nirkabel.
            \end{itemize}
        \end{column}
        \begin{column}{0.48\textwidth}
            \centering
            \bookfig[1.0\linewidth]{02}{penyadapan-akustik.png}
            \figsumber{Munir, \textit{IF4020 Kriptografi}, ITB.}
        \end{column}
    \end{columns}
\end{frame}

\begin{frame}{Serangan Aktif dan \textit{Man-in-the-Middle}}
    \begin{columns}[T,onlytextwidth]
        \begin{column}{0.54\textwidth}
            \slidejudul{Mengintervensi komunikasi}
            \begin{itemize}
                \item Menghapus sebagian cipherteks.
                \item Mengubah cipherteks.
                \item Menyisipkan potongan cipherteks palsu.
                \item Memutar ulang pesan lama.
                \item Mengubah informasi yang tersimpan.
            \end{itemize}
            \begin{alertblock}{Contoh}
                \textit{Man-in-the-middle}: koneksi dialihkan melalui
                penyerang.
            \end{alertblock}
        \end{column}
        \begin{column}{0.44\textwidth}
            \centering
            \bookfig[0.98\linewidth]{02}{serangan-mitm.png}
            \figsumber{Munir, \textit{IF4020 Kriptografi}, ITB.}
        \end{column}
    \end{columns}
    \vspace{0.05cm}
    \begin{block}{Yang membuat MITM berbahaya}
        Kedua pihak merasa berkomunikasi langsung dengan pihak yang benar.
    \end{block}
\end{frame}

\begin{frame}{Serangan Pasif dan Aktif: Perbandingan}
    \footnotesize
    \begin{tabularx}{\linewidth}{@{}>{\raggedright\arraybackslash}p{2.4cm}
        >{\raggedright\arraybackslash}X
        >{\raggedright\arraybackslash}X@{}}
        \toprule
        \textbf{Aspek} & \textbf{Serangan pasif} & \textbf{Serangan aktif} \\
        \midrule
        Keterlibatan & Tidak terlibat; hanya menyadap. & Mengintervensi dan
        memengaruhi sistem. \\
        \addlinespace
        Bentuk & Intersepsi dan analisis lalu lintas. & Menghapus, mengubah,
        atau menyisipkan cipherteks; memutar ulang pesan; mengubah informasi
        tersimpan; MITM. \\
        \addlinespace
        Data & Tidak diubah. & Diubah, ditambah, atau dihapus. \\
        \bottomrule
    \end{tabularx}
    \vspace{0.2cm}
    \begin{alertblock}{Poin}
        Klasifikasi ini berasal dari keterlibatan penyerang dalam komunikasi.
    \end{alertblock}
\end{frame}

% =====================================================================
\section{Keamanan Komputasi dan Serangan Saluran Samping}
% =====================================================================

\begin{frame}{Keamanan Komputasi}
    \begin{block}{Aman secara komputasi}
        Sebuah algoritma aman secara komputasi bila memenuhi tiga syarat.
    \end{block}
    \vspace{0.15cm}
    \begin{enumerate}
        \item Persamaan matematika yang menggambarkan operasi algoritma sangat
              kompleks sehingga tidak mungkin dipecahkan secara analitik.
        \item Biaya untuk memecahkan cipherteks melampaui nilai informasi
              yang terkandung di dalamnya.
        \item Waktu yang diperlukan untuk memecahkan cipherteks melampaui
              lamanya informasi harus dijaga kerahasiaannya.
    \end{enumerate}
    \vspace{0.1cm}
    \begin{alertblock}{Konteks}
        Ketiga syarat bergantung pada keadaan. Algoritma yang sama dapat
        dianggap aman pada satu situasi dan tidak aman pada situasi lain.
    \end{alertblock}
\end{frame}

\begin{frame}{Serangan Saluran Samping}
    \footnotesize
    Serangan saluran samping tidak menyerang matematika algoritmanya, tetapi
    memanfaatkan informasi yang bocor ketika algoritma dijalankan pada
    perangkat fisik.
    \vspace{0.1cm}
    \begin{tabularx}{\linewidth}{@{}>{\raggedright\arraybackslash}p{3.0cm}
        >{\raggedright\arraybackslash}X@{}}
        \toprule
        \textbf{Jenis} & \textbf{Informasi yang dimanfaatkan} \\
        \midrule
        Pewaktuan (\textit{timing}) & Perbedaan lama eksekusi operasi. \\
        \addlinespace
        Daya (\textit{power}) & Pola konsumsi daya perangkat; ragamnya SPA
        dan DPA. \\
        \addlinespace
        Elektromagnetik & Radiasi elektromagnetik yang dihasilkan perangkat. \\
        \addlinespace
        Penyuntikan kesalahan & Keluaran yang menyimpang akibat kesalahan
        buatan; contohnya DFA. \\
        \bottomrule
    \end{tabularx}
    \vspace{0.15cm}
    \begin{alertblock}{Poin}
        Algoritma yang kuat secara matematis tetap dapat bocor melalui cara
        pelaksanaannya pada perangkat fisik.
    \end{alertblock}
\end{frame}

% ------------------------------------------------------------------ penutup
\begin{frame}{Yang Harus Diingat}
    \footnotesize
    \begin{block}{Yang harus diingat}
        \begin{itemize}
            \item Serangan adalah usaha menemukan kunci atau plainteks;
                  Prinsip Kerckhoffs menegaskan algoritma publik dan hanya
                  kunci yang rahasia.
            \item Berdasarkan informasi: \textit{ciphertext-only},
                  \textit{known-plaintext}, \textit{chosen-plaintext},
                  \textit{adaptive chosen-plaintext}, dan
                  \textit{chosen-ciphertext}.
            \item Berdasarkan tujuan: \textit{key-recovery},
                  \textit{plaintext-recovery}, \textit{distinguishing},
                  \textit{forgery}, dan \textit{replay}.
            \item Serangan pasif menyadap; serangan aktif mengintervensi,
                  termasuk \textit{man-in-the-middle}.
            \item Serangan saluran samping: pewaktuan, daya, elektromagnetik,
                  dan penyuntikan kesalahan.
        \end{itemize}
    \end{block}
    \vspace{0.35cm}
    \begin{exampleblock}{Pertemuan berikutnya}
        \textbf{Minggu 3 --- Landasan matematika}
    \end{exampleblock}
\end{frame}

\end{document}
